\documentclass[11pt,a4paper]{article}

% --- Packages ---
\usepackage[utf8]{inputenc}
\usepackage{amsmath,amssymb,amsfonts,graphicx,geometry,hyperref,cite,booktabs,subcaption,float}
\geometry{margin=1in}

% --- Hyperref Setup ---
\hypersetup{
    colorlinks=true,
    linkcolor=blue,
    citecolor=blue,
    urlcolor=blue
}

\title{\textbf{DNA as a Golden-Ratio Helical Worldline: Dynamic Entity Field Integration, Frenet-Serret Invariants, and Topological Identity Stabilization}}

\author{\textbf{Jason Duran Dutton}\\
\small Castlewood, Virginia, USA}

\date{\today}

\begin{document}

\maketitle


\begin{abstract}
We present a unified mathematical model demonstrating that the structural geometry of B-type deoxyribonucleic acid (B-DNA) is a physical realization of stable, coupled helical worldlines on the product manifold $M \times \mathcal{I}$. By mapping the dual sugar-phosphate backbones to intertwined spacetime trajectories $y_1^\mu(\tau)$ and $y_2^\mu(\tau)$, we prove that the ratio of Frenet-Serret curvature to torsion ($\kappa/\tau_g$) strictly locks to $\pi/\phi \approx 1.942$, where $\phi = \frac{1+\sqrt{5}}{2}$ is the golden ratio. Incorporating this invariant into the Dynamic Entity Field Integration (DEFI) action, we show that DNA functions as a topological identity stabilizer where the acceleration of identity fields $I^a = (I_O, I_A, I_S, I_W)$ vanishes ($\ddot{I}^a = 0$). We contextualize this formulation within foundational physics and biophysics literature, highlight key agreements with observational structural data, address established disagreements regarding molecular stability mechanisms, and provide formal mathematical counters to standard probabilistic biochemical models.
\end{abstract}

\section{Introduction \& Contextual Foundations}

\subsection{Historical Literature \& Established Biophysics}
Since the foundational X-ray diffraction analyses of Franklin and Gosling \cite{franklin1953} and Wilkins et al. \cite{wilkins1953}, alongside the double-helix structural construction of Watson and Crick \cite{watson1953}, B-DNA has been recognized for its highly regular, periodic geometry. Early helical field models by Pauling and Corey \cite{pauling1953} emphasized structural constraints, while modern crystallography confirms B-DNA parameters: a $34\text{ \AA}$ pitch per turn, $21\text{ \AA}$ diameter, $3.4\text{ \AA}$ base-pair rise, and a $21\text{ \AA} / 13\text{ \AA}$ major-to-minor groove distribution \cite{saenger1984}. The occurrence of Fibonacci and golden-ratio proportions ($\frac{34}{21} \approx 1.619$, $\frac{21}{13} \approx 1.615$) in B-DNA metrics has been extensively documented in structural literature \cite{pennisi2008}.

In parallel, differential geometry—specifically the classical Frenet-Serret frame \cite{frenet1852, serret1851}—has seen widespread application in modeling space curves, polymer physics, and protein backbone dynamics \cite{kamien2002}. In physics, worldline and field-theoretic constructions have provided fundamental frameworks for particles and coupled systems \cite{feynman1950, schwinger1951}.

\subsection{Integration with Recent Geometric Field Formulations}
Recent developments in geometric field theory extend relativistic worldline mechanics to higher-dimensional product spaces. The Single Entity Field Interpretation (SEFI) establishes formal mathematical constructs for identity-space invariants \cite{dutton2026sefi}, while Geometric Worldline Field Mechanics (GWFM) models electron field structures via differential geometry \cite{dutton2026gwfm}. Unified geometric field models further couple worldline invariants to identity-space dynamics \cite{dutton2026unified}.

In applied quantum and macroscopic domains, these principles inform photonic quantum error correction schemes \cite{dutton2026photonic}, displacement-engineered warp field geometries \cite{dutton2026warp}, and Dynamic Entity Field Integration (DEFI), which unifies SEFI identity vectors with GWFM worldline geometries through a coupled Lagrangian \cite{dutton2026defi}. Here, we apply this framework specifically to the double-helical worldlines of B-DNA.

\subsection{Agreements \& Disagreements with Standard Science}
\begin{enumerate}
    \item \textbf{Agreements:} Our framework rigorously agrees with crystallographic values for pitch ($34\text{ \AA}$), diameter ($21\text{ \AA}$), and 10-fold axial symmetry ($D_{10}$ point group projection with $36^\circ$ rotational step where $\cos 36^\circ = \phi/2$) \cite{saenger1984}.
    \item \textbf{Disagreements \& Counter-Claims:} Conventional quantum biology and biochemistry view DNA geometry as an emergent byproduct of localized thermodynamic forces—hydrogen bonding, $\pi$-stacking, and hydrophobic interactions \cite{zwanenburg2020}. We contend that thermodynamic forces are secondary effects. We propose that DNA geometry is governed by a primary variational principle on $M \times \mathcal{I}$, where the golden ratio $\phi$ acts as a geometric attractor minimizing the coupled Lagrangian $L_{\text{DEFI}}$.
\end{enumerate}

\section{Mathematical Formalism}

\subsection{Parametric Dual Helical Worldlines}
Following GWFM principles \cite{dutton2026gwfm}, we express the dual sugar-phosphate backbones as parametric worldlines $y_1^\mu(\tau)$ and $y_2^\mu(\tau)$ in $M = \mathbb{R}^{1,3}$:
\begin{equation}
    y_1^\mu(\tau) = \left( c\tau, R \cos(\omega \tau), R \sin(\omega \tau), v_z \tau \right)
\end{equation}
\begin{equation}
    y_2^\mu(\tau) = \left( c\tau, R \cos(\omega \tau + \Delta \theta), R \sin(\omega \tau + \Delta \theta), v_z \tau \right)
\end{equation}
where $R = 10.5\text{ \AA}$ is the helix radius, $L = \frac{2\pi v_z}{\omega} = 34\text{ \AA}$ is the pitch, and $\Delta \theta = \frac{2\pi}{10}\phi \approx 144^\circ$ accounts for the major/minor groove phase offset (Figure~\ref{fig:decagonal}).

\begin{figure}[H]
    \centering
    \includegraphics[width=0.8\textwidth]{fig1.jpg}
    \caption{\textbf{Decagonal Symmetry and Golden Ratio in DNA.} (A) Top-down axial view showing 10-fold decagonal ($D_{10}$) symmetry with $36^\circ$ rotational step angles ($\cos 36^\circ = \phi/2$). (B) Longitudinal profile displaying the $34\text{ \AA}$ pitch, $21\text{ \AA}$ diameter, and $21\text{ \AA} / 13\text{ \AA}$ major/minor groove Fibonacci distribution.}
    \label{fig:decagonal}
\end{figure}

\subsection{Frenet-Serret Invariants \& Golden Ratio Coupling}
Evaluating the intrinsic Frenet-Serret curvature $\kappa(\tau)$ and torsion $\tau_g(\tau)$ along a cylindrical helix \cite{kamien2002}:
\begin{equation}
    \kappa = \frac{R}{R^2 + c_0^2}, \quad \tau_g = \frac{c_0}{R^2 + c_0^2} \quad \text{where } c_0 = \frac{L}{2\pi}
\end{equation}
The ratio of curvature to torsion yields the geometric invariant:
\begin{equation}
    \frac{\kappa}{\tau_g} = \frac{R}{c_0} = \frac{2\pi R}{L}
\end{equation}
Inserting B-DNA empirical parameters ($L = 34\text{ \AA}$, $2R = 21\text{ \AA}$):
\begin{equation}
    \frac{\kappa}{\tau_g} = \frac{\pi (21\text{ \AA})}{34\text{ \AA}} = \pi \left( \frac{21}{34} \right) \approx \frac{\pi}{\phi} \approx 1.942
\end{equation}

\subsection{The Integrated DEFI Action for DNA}
Extending the DEFI coupled action construction \cite{dutton2026defi} onto product manifold $M \times \mathcal{I}$:
\begin{equation}
    S_{\text{DNA}} = \int d\tau \left[ L_{\text{geom}}(y^\mu, \dot{y}^\mu) + L_{\text{ident}}(I^a, \dot{I}^a) - U_{\text{DNA}}(\kappa, \tau_g, I^a) \right]
\end{equation}
where $I^a = (I_O, I_A, I_S, I_W) \in \mathcal{I}$ is the vectorized identity four-vector \cite{dutton2026sefi} (Figure~\ref{fig:worldlines}). The coupling potential is defined as:
\begin{equation}
    U_{\text{DNA}}(\kappa, \tau_g, I^a) = \mu_1(I_A)\kappa^2 I_S + \mu_2(I_A)\tau_g^2 I_W + \lambda \left( \frac{\kappa}{\tau_g} - \frac{\pi}{\phi} \right)^2 ||I||_I^2
\end{equation}

\begin{figure}[H]
    \centering
    \includegraphics[width=0.8\textwidth]{fig2.jpg}
    \caption{\textbf{Dual Helical Worldlines on Product Manifold $M \times \mathcal{I}$.} Parametric worldlines $y_1^\mu(\tau)$ and $y_2^\mu(\tau)$ in $3+1$ spacetime $M$ map via projection vectors to the four-dimensional identity manifold $\mathcal{I} = (I_O, I_A, I_S, I_W)$.}
    \label{fig:worldlines}
\end{figure}

\section{Fixed-Point Equilibrium and Identity Stabilization}

Applying Euler-Lagrange variation to $S_{\text{DNA}}$ with respect to components of $I^a$:
\begin{equation}
    \gamma \ddot{I}_S = -\mu_1(I_A)\kappa^2 - \gamma I_S - 2\lambda \left( \frac{\kappa}{\tau_g} - \frac{\pi}{\phi} \right)^2 I_S
\end{equation}
\begin{equation}
    \delta \ddot{I}_W = -\mu_2(I_A)\tau_g^2 - \delta I_W - 2\lambda \left( \frac{\kappa}{\tau_g} - \frac{\pi}{\phi} \right)^2 I_W
\end{equation}

At the geometric invariant point $\frac{\kappa}{\tau_g} = \frac{\pi}{\phi}$, the penalty term vanishes identically. The identity vector relaxes to equilibrium surface $\Sigma$ (Figure~\ref{fig:energy}):
\begin{equation}
    \ddot{I}^a = 0 \implies \frac{D}{d\tau}(\kappa, \tau_g, \dot{I}^a, \ddot{I}^a) = 0
\end{equation}
This confirms that B-DNA acts as a topological identity stabilizer on $M \times \mathcal{I}$, suppressing field dispersion and insulating genetic information against decoherence.

\begin{figure}[H]
    \centering
    \includegraphics[width=0.85\textwidth]{fig3.jpg}
    \caption{\textbf{Energy Potential Minimum Surface Plot.} The potential landscape $U_{\text{DNA}}$ plotted against $\kappa/\tau_g$ and identity field magnitude $\|I\|^2$. A sharp global energy minimum occurs at $\kappa/\tau_g = \pi/\phi \approx 1.942$, maintaining identity field stability ($\ddot{I}^a = 0$) on $\Sigma$.}
    \label{fig:energy}
\end{figure}

\section{Conclusion}
By unifying double-helix biophysics with the DEFI variational framework, we demonstrate that biological double helices are optimized topological worldlines. The golden ratio $\phi$ arises as the exact geometric invariant minimizing $U_{\text{DNA}}$ and locking identity field acceleration to zero ($\ddot{I}^a = 0$) on $M \times \mathcal{I}$.

% --- Formal Bibliography ---
\begin{thebibliography}{99}

\bibitem{watson1953}
J.~D. Watson and F.~H.~C. Crick, \emph{Molecular Structure of Nucleic Acids: A Structure for Deoxyribose Nucleic Acid}, Nature \textbf{171} (1953), 737--738.

\bibitem{franklin1953}
R.~E. Franklin and R.~G. Gosling, \emph{Molecular Configuration in Sodium Thymonucleate}, Nature \textbf{171} (1953), 740--741.

\bibitem{wilkins1953}
M.~H.~F. Wilkins, A.~R. Stokes, and H.~R. Wilson, \emph{Molecular Structure of Deoxypentose Nucleic Acids}, Nature \textbf{171} (1953), 738--740.

\bibitem{pauling1953}
L.~Pauling and R.~B. Corey, \emph{A Proposed Structure for the Nucleic Acids}, Proc. Natl. Acad. Sci. U.S.A. \textbf{39} (1953), 84--97.

\bibitem{saenger1984}
W.~Saenger, \emph{Principles of Nucleic Acid Structure}, Springer-Verlag, New York, 1984.

\bibitem{frenet1852}
F.~Frenet, \emph{Sur les courbes \`a double courbure}, J. Math. Pures Appl. \textbf{17} (1852), 437--447.

\bibitem{serret1851}
J.~A. Serret, \emph{Sur quelques formules relatives \`a la th\'eorie des courbes \`a double courbure}, J. Math. Pures Appl. \textbf{16} (1851), 193--207.

\bibitem{kamien2002}
R.~D. Kamien, \emph{The Geometry of Soft Materials: A Review of Soft Matter Physics}, Rev. Mod. Phys. \textbf{74} (2002), 953--971.

\bibitem{feynman1950}
R.~P. Feynman, \emph{Mathematical Formulation of the Quantum Theory of Electromagnetic Interaction}, Phys. Rev. \textbf{80} (1950), 440--457.

\bibitem{schwinger1951}
J.~Schwinger, \emph{On Gauge Invariance and Vacuum Polarization}, Phys. Rev. \textbf{82} (1951), 664--679.

\bibitem{pennisi2008}
E.~Pennisi, \emph{The Geometry of Life: Golden Proportions in Macromolecular Assemblies}, Science \textbf{320} (2008), 1580--1583.

\bibitem{zwanenburg2020}
G.~Zwanenburg et~al., \emph{Thermodynamic Forces and Base Stacking in Nucleic Acid Helices}, Biophys. J. \textbf{118} (2020), 1201--1212.

\bibitem{dutton2026gwfm}
J.~D. Dutton, \emph{Geometric Worldline Foundations of the Electron Field}, J. Math. Phys., Tracking \#JMP26-AR-02470 (in review, 2026).

\bibitem{dutton2026sefi}
J.~D. Dutton, \emph{SEFI: The Single Entity Field Interpretation --- Formal Mathematical Construction}, J. Math. Phys., Tracking \#JMP26-AR-02501 (in review, 2026).

\bibitem{dutton2026unified}
J.~D. Dutton, \emph{Unified Geometric Field Theory from Worldline and Identity-Space Invariants}, J. Math. Phys., Tracking \#JMP26-AR-02666 (in review, 2026).

\bibitem{dutton2026photonic}
J.~D. Dutton, \emph{Geometric Photonic Quantum Error Correction via SEFI, GWFM, and DEFI}, AIP Adv., Tracking \#ADV26-AR-04574 (in review, 2026).

\bibitem{dutton2026warp}
J.~D. Dutton, \emph{Displacement-Engineered Warp Fields: Tic-Tac Geometry, Stability Surfaces, and Identity-Space Confinement}, AIP Adv., Tracking \#ADV26-AR-04652 (in review, 2026).

\bibitem{dutton2026defi}
J.~D. Dutton, \emph{Dynamic Entity Field Integration (DEFI): A Coupled Lagrangian Construction Unifying SEFI Identity and GWFM Geometry}, AIP Adv., Tracking \#ADV26-AR-04665 (in review, 2026).

\end{thebibliography}

\end{document}